\chapter{Discussion, Limitations and Future Work}\label{ch:discussion}

This chapter interprets the results, states the prototype's limitations, and identifies priorities for
completing, deploying, and extending the system.

% ============================================================================
\section{Discussion}\label{sec:disc-discussion}

DUET adds automatic in-band substitution detection through additional label-bound lookups rather than
a new proof family, so it preserves the lookup proof's asymptotic size and verification cost. The
guarantee depends on domain-separated labels, proof-verified absence, warning persistence, and
transition validation. Under \hyperref[sc:threat-a3]{A3}, divergent conversation secrets produce different labels, while the
global VRF-bound identity entry supplies an independent binding.

The evaluation supports this separation of costs. Identity proofs remain roughly comparable with the
matched CONIKS proof, detection paths reuse raw-label proofs, and distant-head consistency replaces
a proof per missed epoch with one proof. These measurements characterize component costs rather than
production throughput.

The construction applies beyond Signal wherever a messenger supplies a semi-trusted directory and a
conversation secret, with continuous mode additionally needing public ratchet keys and secret epoch
inputs. Session mode has the smaller integration and write surface, while continuous mode
authenticates each ratchet epoch at a higher write rate. Both remain bounded by trust on first use,
endpoint security, and the need for an independent vantage to expose a consistent fork.

% ============================================================================
\section{Limitations}\label{sec:disc-limitations}

The guarantee begins only after a baseline has been pinned, so it cannot establish whether the
initial binding was authentic. It covers
pairwise, single-device conversations rather than groups, MLS, Sender Keys, or complete multi-device
warning propagation. Endpoint compromise remains outside any directory guarantee.

Transition auditing requires an honest auditor, and \hyperref[sc:threat-a4]{A4} additionally requires a valid signed head from
an independent vantage. A client alone rejects a head that does not extend its pinned baseline, but
it cannot expose two internally consistent branches extending a common prefix. Production deployment
also needs authenticated epoch disclosures, freshness and failure policies, and a real gossip or
second-vantage source.

The prototype provides no query privacy from the operator, write-rate limit, published-head timeout,
or gossip-based proof-omission detection. A malicious peer can publish a valid but false warning,
since the payload MAC authenticates the peer rather than the truth of its report, so attributable
warning records and abuse policy remain necessary.

The implementation is single-machine and in-memory, and epoch commit deep-clones the current tree into
the committed snapshot. The Signal client uses a deterministic conversation-secret stand-in and a
seeded ratchet simulation, which exercise the code
paths but do not establish label unpredictability or active-attacker security for live Signal state.
The key-fetch hook reports a rejecting verdict without blocking session establishment. The simulated
continuous mode supplies the raw Diffie--Hellman output where ACKA's Signal mapping uses the derived
chain key, and its synchronous state does not model skipped ratchet epochs. A faithful binding must
obtain the live chain key and support explicit epoch indices and bounded skipped state.

The security argument is an informal proof sketch supported by definitions, the same-head lemma, and
the adversarial suite. The suite does not model separate first-contact sessions with each party or a
fabricated pre-key bundle kept consistent with a malicious directory entry. Timings measure local
components rather than a complete networked monitor or audit, and the committed corpus contains
aggregate statistics rather than raw Criterion samples. M11 repeats the Rust lookup-verification
component rather than timing the TypeScript key-fetch path, and M12 measures local signature and
consistency work rather than network transfer, epoch-delta retrieval, or full transition replay.
These distinctions prevent the component evidence from being read as an end-to-end deployment
measurement.

% ============================================================================
\section{Future Work}\label{sec:disc-future}

\subsection{Completing the Protocol and Security Argument}

The first implementation priority is to obtain the live X3DH or PQXDH conversation secret and the
live Double Ratchet chain key from libsignal. Replacing the stand-ins changes the inputs to the
existing label and CEA constructions rather than their downstream proof formats. The live
conversation secret is required for the \hyperref[sc:threat-a3]{A3} label-unpredictability claim. The
raw-label interface must also authenticate the party authorized to make the first write, preventing a
malicious client from occupying an empty conversation slot before its legitimate owner. Later
replacement and deletion remain governed by transition auditing rather than first-write
authorization. The tested persistence component provides evidence of later removal, and its
implementation status is identified in \cref{tab:impl-modules}. A formal analysis should define a
combined detection game covering the KDF, VRF, signed heads, warning authentication, and append-only
verification, alongside symbolic analysis and separate verification of the client proof checker.

\subsection{Groups, Devices, and Directory Privacy}

Group support requires epoch-scoped labels and a
warning policy for joins, removals, and every current device, potentially using domain-separated MLS
exporter secrets~\cite{rfc9420}. ELEKTRA-style device cross-signing could strengthen multi-device
continuity~\cite{len2023elektra}, alongside device-scoped labels, warning propagation, and reset
handling. OPRF-based identity blinding could hide the raw lookup identity, but private retrieval and
network anonymity would still be needed to conceal the fetched entry and the requester, and the
design would have to define how the OPRF role is separated or distributed~\cite{rfc9497}.

\subsection{Deployment and Cryptographic Scale}

Persistent copy-on-write nodes should replace full
snapshot cloning. With fixed depth $d=256$, copying the paths for $k$ modified labels is $O(kd)$ and
therefore $O(k)$, independent of $N$. The proof format would remain unchanged. Building on Signal's
existing key-transparency service~\cite{signal-kt-server} would additionally
require immutable raw writes, authenticated ownership, transition disclosures, client hooks, and an
independent auditor or gossip federation with incentives for independent operation. PQXDH secrets
already fit the label KDF, but post-quantum operation also needs replacements for the ECVRF and
signed-head signatures~\cite{pqxdh}. This thesis does not select or implement a post-quantum VRF.
Even with PQXDH-derived labels, end-to-end post-quantum security would still depend on the messenger's
authentication and ratchet mechanisms. Concurrency, serialization, reclamation, and memory use would
require measurement after structural sharing.

\subsection{Evaluation, Operations, and Usability}

A usability study should determine whether cryptographic detection leads to an appropriate user
response. It should test interpretation of the red warning and amber notice, accessibility,
escalation, and alarm fatigue, continuing the ceremony research that motivated the
design~\cite{vaziripour2017}. Production warnings
should carry the scenario, epoch, signed heads, and proofs needed for operator or auditor review
without exposing unnecessary metadata. An operational policy must define when to reject a proof,
request rotation, or escalate without turning warning records into a new metadata leak. Further
experiments should use persistent storage, concurrent requests, wide-area delay, varied epoch and
catch-up policies, mobile CPU and energy measurements, and archived raw samples. The same approach
may also be evaluated in clinical, emergency-response, and industrial key-distribution settings, each
with its own identity model, device-enrolment and revocation process, availability requirements, and
regulatory controls.
